% Filename: lista01.tex
% 
% This code is part of 'Solutions for MT402, Matrizes'
% 
% Description: This file corresponds to the solutions of homework sheet 01.
% 
% Created: 01.03.12 04:00:00 PM
% Last Change: 04.06.12 05:11:05 PM
% 
% Authors:
% - Raniere Silva (2012): initial version
% 
% Copyright (c) 2012 Raniere Silva <r.gaia.cs@gmail.com>
% 
% This work is licensed under the Creative Commons Attribution-ShareAlike 3.0 Unported License. To view a copy of this license, visit http://creativecommons.org/licenses/by-sa/3.0/ or send a letter to Creative Commons, 444 Castro Street, Suite 900, Mountain View, California, 94041, USA.
%
% This work is distributed in the hope that it will be useful, but WITHOUT ANY WARRANTY; without even the implied warranty of MERCHANTABILITY or FITNESS FOR A PARTICULAR PURPOSE.
%
\documentclass[a4paper,12pt, leqno, answers]{exam}
% Customiza\c{c}\~{a}o da classe exam
\newcommand{\mycheader}{Lista 01 - Revis\~{a}o de \'{A}lgebra Matricial I}
\header{MT402}{\mycheader}{\thepage/\numpages}
\headrule
\footer{Dispon\'{i}vel em \\\url{https://github.com/r-gaia-cs/solucoes_lista_matrizes}
}{}{Reportar erros para \\\href{mailto:r.gaia.cs@gmail.com}{r.gaia.cs@gmail.com}
}
\footrule 
\pagestyle{headandfoot}
\renewcommand{\solutiontitle}{\noindent\textbf{Solu\c{c}\~{a}o:}\enspace}
\SolutionEmphasis{\slshape}
\unframedsolutions
\pointname{}

% Filename: paper_size.tex
% 
% This code is part of 'Solutions for MT402, Matrizes'
% 
% Description: This file corresponds to the paper size output.
% 
% Created: 08.05.12 09:46:55 PM
% Last Change: 04.06.12 10:42:01 PM
% 
% Authors:
% - Raniere Silva (2012): initial version
% 
% Copyright (c) 2012 Raniere Silva <r.gaia.cs@gmail.com>
% 
% This work is licensed under the Creative Commons Attribution-ShareAlike 3.0 Unported License. To view a copy of this license, visit http://creativecommons.org/licenses/by-sa/3.0/ or send a letter to Creative Commons, 444 Castro Street, Suite 900, Mountain View, California, 94041, USA.
%
% This work is distributed in the hope that it will be useful, but WITHOUT ANY WARRANTY; without even the implied warranty of MERCHANTABILITY or FITNESS FOR A PARTICULAR PURPOSE.
%
% Para impress\~{a}o
\usepackage[top=3cm, bottom=3cm, left=2cm, right=2cm]{geometry}

% Para ereaders (Kindle, Nook, Kobo, ...) and tablets (iPad, GalaxyTab, ...)
% \usepackage[papersize={180mm,240mm},margin=2mm]{geometry}
% \sloppy


% Pacotes
\usepackage[utf8]{inputenc}
\usepackage[T1]{fontenc}
\usepackage[brazil]{babel}
\usepackage{amsmath}
\usepackage{amsfonts}
\usepackage{amssymb}
\usepackage{hyperref}
\usepackage{graphicx}

% Customiza\c{c}\~{a}o do pacote amsmath
\allowdisplaybreaks[4]

% Novos ambientes
% \newenvironment{fwsolution}{\begin{EnvFullwidth}\begin{TheSolution}}{\end{TheSolution}\end{EnvFullwidth}}

% Novos comandos
% \newcommand{\mdot}{\text{\LARGE $\boldsymbol{\cdot}$}}
\newcommand{\mdot}{\bullet}

\begin{document}
%cover
\thispagestyle{empty}
% Filename: cover.tex
% 
% This code is part of 'Solutions for MT402, Matrizes'
% 
% Description: This file corresponds to the cover.
% 
% Created: 30.05.12 04:40:25 PM
% Last Change: 14.07.12 10:49:33 AM
% 
% Authors:
% - Raniere Silva (2012): initial version
% 
% Copyright (c) 2012 Raniere Silva <r.gaia.cs@gmail.com>
% 
% This work is licensed under the Creative Commons Attribution-ShareAlike 3.0 Unported License. To view a copy of this license, visit http://creativecommons.org/licenses/by-sa/3.0/ or send a letter to Creative Commons, 444 Castro Street, Suite 900, Mountain View, California, 94041, USA.
%
% This work is distributed in the hope that it will be useful, but WITHOUT ANY WARRANTY; without even the implied warranty of MERCHANTABILITY or FITNESS FOR A PARTICULAR PURPOSE.
%
\begin{center}
    \LARGE{Solu\c{c}\~{o}es para MT402, Matrizes}
    
    \Large{\mycheader}
\end{center}
\vspace{.5\textheight}

\begin{tabular}{|p{.9\textwidth}|}
\hline
Este trabalho foi licenciado com a Licen\c{c}a Creative Commons Atribui\c{c}\~{a}o - CompartilhaIgual 3.0 N\~{a}o Adaptada. Para ver uma c\'{o}pia desta licen\c{c}a, visite \url{http://creativecommons.org/licenses/by-sa/3.0/} ou envie um pedido por carta para Creative Commons, 444 Castro Street, Suite 900, Mountain View, California, 94041, USA.
\begin{center}
\includegraphics[scale=1]{cc-by-sa.png}
\end{center}
Este trabalho encontra-se dispon\'{i}vel em \url{https://github.com/r-gaia-cs/solucoes_lista_matrizes}
 e atualmente precisa de um mantenedor, se tiver interessse entre em contato com Raniere Silva (\href{mailto:r.gaia.cs@gmail.com}{r.gaia.cs@gmail.com}
).

Este trabalho \'{e} distribuido na esperança que possa ser \'{u}til, mas SEM NENHUMA GARANTIA; sem uma garantia implicita de ADEQUA\c{C}\~{A}O a qualquer MERCADO ou APLICA\c{C}\~{A}O EM PARTICULAR.
\\ \hline
\end{tabular}



\newpage
\setcounter{page}{1}
\begin{questions}
    \question Algumas defini\c{c}\~{o}es:
    \begin{description}
        \item[matriz conjugada:] a conjugada de uma matriz $A = (a_{ij})$ \'{e} denotada por $\overline{A}$ e \'{e} definida por $\overline{A} = (\overline{a}_{ij})$.
        \item[transposta conjugada:] \'{e} definida por $\overline{A^t}$ e \'{e} denotada por $A^*$ (livro do Trefethen\nocite{Trefethen:1997:numerical} e Mayer\nocite{Meyer:2000:matrix}) ou $A^H$ (livro do Golub\nocite{Golub:1996:matrix}).
        \item[sim\'{e}trica] quando $A = A^t$.
        \item[anti-sim\'{e}trica] quando $A = - A^t$.
        \item[Hermitiana] quando $A = A^*$.
        \item[anti-Hermitiana] quando $A = - A^*$.
    \end{description}

    D\^{e} um exemplo de cada matriz considerando o espa\c{c}o das matrizes $\mathbb{C}^{4 \times 4}$.
    \begin{parts}
        \part sim\'{e}trica
        \begin{solution}
            \[
            \begin{bmatrix}
                1 & i & 1 & 0 \\
                i & 0 & -i & 1 \\
                1 & -i & -1 & 0 \\
                0 & 1 & 0 & i
            \end{bmatrix}
            \]
        \end{solution}

        \part anti-sim\'{e}trica
        \begin{solution}
            \[
            \begin{bmatrix}
                0 & -i & -1 & 0 \\
                i & 0 & i & -1 \\
                1 & -i & 0 & 0 \\
                0 & 1 & 0 & 0
            \end{bmatrix}
            \]
        \end{solution}

        \part Hermitiana
        \begin{solution}
            \[
            \begin{bmatrix}
                1 & -i & -1 & 0 \\
                i & 0 & i & -1 \\
                -1 & -i & 5 & 0 \\
                0 & -1 & 0 & -2
            \end{bmatrix}
            \]
        \end{solution}

        \part anti-Hermitiana
        \begin{solution}
            \[
            \begin{bmatrix}
                0 & -i & -1 & 0 \\
                -i & 2i & i & 1 \\
                1 & i & -3i & 0 \\
                0 & -1 & 0 & 0
            \end{bmatrix}
            \]
        \end{solution}

        \part Hermitiana mas n\~{a}o sim\'{e}trica
        \begin{solution}
            \[
            \begin{bmatrix}
                1 & -i & -1 & 0 \\
                i & 0 & i & -1 \\
                -1 & -i & 5 & 0 \\
                0 & -1 & 0 & -2
            \end{bmatrix}
            \]
        \end{solution}

        \part sim\'{e}trica mas n\~{a}o Hermitiana
        \begin{solution}
            \[
            \begin{bmatrix}
                1 & i & 1 & 0 \\
                i & 0 & -i & 1 \\
                1 & -i & -1 & 0 \\
                0 & 1 & 0 & i
            \end{bmatrix}
            \]
        \end{solution}

        \part sim\'{e}trica e anti-sim\'{e}trica
        \begin{solution}
            \[
            \begin{bmatrix}
                0 & 0 & 0 & 0 \\
                0 & 0 & 0 & 0 \\
                0 & 0 & 0 & 0 \\
                0 & 0 & 0 & 0
            \end{bmatrix}
            \]
        \end{solution}

        \part sim\'{e}trica e anti-Hermitiana
        \begin{solution}
            \[
            \begin{bmatrix}
                0 & i & i & -i \\
                i & 0 & -i & 0 \\
                i & -i & 0 & 0 \\
                -i & 0 & 0 & 0
            \end{bmatrix}
            \]
        \end{solution}
    \end{parts}

    \question Julgue verdadeiro ou falso a afirma\c{c}\~{a}o abaixo. Justifique.
    \begin{quote}
        Se os coeficientes da matriz $A$ s\~{a}o todos reais e $A$ \'{e} sim\'{e}trica, ent\~{a}o A \'{e} tamb\'{e}m Hermitiana.
    \end{quote}
    \begin{solution}
        A afirma\c{c}\~{a}o \'{e} verdadeira. Se $A \in \mathbb{R}^{n \times n}$ ent\~{a}o $A = \overline{A}$ e se $A$ \'{e} sim\'{e}trica ent\~{a}o $A = A^t$ concluimos que $A = \overline{A^t}$, i.e., $A$ \'{e} Hermitiana.
    \end{solution}

    \question[Exerc\'{i}cio 3.2.5, p\'{a}gina 88, do Mayer\nocite{Meyer:2000:matrix}] Demostre que:
    \begin{parts}
        \part se $A$ \'{e} anti-sim\'{e}trica, ent\~{a}o, $a_{ii} = 0$;
        \begin{solution}
            Se $A$ \'{e} anti-sim\'{e}trica ent\~{a}o $A = - A^t$, isso \'{e}, $a_{ij} = - a_{ji}$. Para $i = j$ deve ocorrer $a_{ii} = - a_{ii}$ que \'{e} poss\'{i}vel apenas se $a_{ii} = 0$.
        \end{solution}

        \part se $A$ \'{e} Hermitiana, ent\~{a}o, $a_{ii}$ \'{e} um n\'{u}mero real;
        \begin{solution}
            Se $A$ \'{e} Hermitiana ent\~{a}o $A = A^*$, isso \'{e}, $a_{ij} = \overline{a_{ij}}$. Para $i = j$ deve ocorrer $a_{ii} = \overline{a_{ii}}$ que \'{e} poss\'{i}vel apenas se $a_{ii} \in \mathbb{R}$.
        \end{solution}

        \part se $A$ \'{e} anti-Hermitiana, ent\~{a}o, $a_{ii}$ \'{e} um n\'{u}mero imagin\'{a}rio puro ($a_{ii} = \gamma_i i$), $\gamma_i \in \mathbb{R}$.
        \begin{solution}
            Se $A$ \'{e} anti-Hermitiana ent\~{a}o $A = - A^*$, isso \'{e}, $a_{ij} = - \overline{a_{ij}}$. Para $i = j$ deve ocorrer $a_{ii} = - \overline{a_{ii}}$ que \'{e} poss\'{i}vel apenas se $a_{ii} = \gamma_i i$, $\gamma_i \in \mathbb{R}$.
        \end{solution}
    \end{parts}

    \question[Exerc\'{i}cio 3.2.6, p\'{a}gina 88, do Meyer\nocite{Meyer:2000:matrix}] \hfill
    \begin{parts}
        \part Demonstre que para qualquer matriz $A$, $n \times n$, $(A + A^t)$ \'{e} sim\'{e}trica e $(A - A^t)$ \'{e} anti-sim\'{e}trica.
        \begin{solution}
            Seja $B = A + A^t$, ent\~{a}o $b_{ij} = a_{ij} + a_{ji}$. Logo, $B$ \'{e} sim\'{e}trica pois $b_{ij} = b_{ji}$.

            Seja $C = A - A^t$, ent\~{a}o $c_{ij} = a_{ij} - a_{ji}$. Logo, $C$ \'{e} anti-sim\'{e}trica pois $c_{ij} = - c_{ji}$.
        \end{solution}

        \part Demonstre que \'{e} poss\'{i}vel escrever uma matriz $A$ de uma \'{u}nica maneira como a soma de duas matrizes, uma sim\'{e}trica e outra anti-sim\'{e}trica: $A = P + Q$, onde $P$ \'{e} uma matriz sim\'{e}trica e $Q$ \'{e} anti-sim\'{e}trica.
        \begin{solution}
            Vamos construir as matrizes $P$ e $Q$, $P$ sim\'{e}trica e $Q$ anti-sim\'{e}trica, tal que $A = P + Q$.
            \[
            \begin{bmatrix}
                a_{11} & a_{12} & \ldots & a_{1n} \\
                a_{21} & a_{22} & \ldots & a_{2n} \\
                \vdots & \vdots & \ddots & \vdots \\
                a_{n1} & a_{n2} & \ldots & a_{nn}
            \end{bmatrix} =
            \begin{bmatrix}
                p_{11} & p_{12} & \ldots & p_{1n} \\
                p_{12} & p_{22} & \ldots & p_{2n} \\
                \vdots & \vdots & \ddots & \vdots \\
                p_{1n} & p_{2n} & \ldots & p_{nn}
            \end{bmatrix} +
            \begin{bmatrix}
                0 & q_{12} & \ldots & q_{1n} \\
                -q_{12} & 0 & \ldots & q_{2n} \\
                \vdots & \vdots & \ddots & \vdots \\
                -q_{1n} & -q_{2n} & \ldots & 0
            \end{bmatrix}
            \]

            Para determinar $P$ e $Q$ devemos resolver o sistema linear dado por
            \[
            \left\{
            \begin{aligned}
                p_{ii} &= a_{ii}, \ i = 1, 2, \ldots, n \\
                p_{ij} + q_{ij} &= a_{ij}, \ \forall 1 \leq i < j \leq n \\
                p_{ij} - q_{ij} &= a_{ji}, \ \forall 1 \leq i < j \leq n
            \end{aligned}
            \right. ,
            \]
            cuja solu\c{c}\~{a}o \'{e}
            \[
            \left\{
            \begin{aligned}
                p_{ii} &= a_{ii}, \ i = 1, 2, \ldots, n \\
                p_{ij} &= \frac{a_{ij} + a_{ji}}{2}, \ \forall 1 \leq i < j \leq n \\
                q_{ij} &= \frac{a_{ij} - a_{ji}}{2}, \ \forall 1 \leq i < j \leq n
            \end{aligned}
            \right. .
            \]

            Pela solu\c{c}\~{a}o do sistema linear verifica-se que \'{e} poss\'{i}vel escrever uma matriz $A$ de uma \'{u}nica maneira como a soma de duas matrizes, uma sim\'{e}trica e outra anti-sim\'{e}trica.
        \end{solution}
    \end{parts}

    \question Considere as matrizes $A : m \times p$ e $B : p \times n$.
    \begin{parts}
        \part Demonstre que cada coluna da matriz $C = A B$ \'{e} uma combina\c{c}\~{a}o linear das colunas de $A$. Quais s\~{a}o os coeficientes desta combina\c{c}\~{a}o linear?
        \begin{solution}
            Pela defini\c{c}\~{a}o de produto matricial temos que
            \[
            C_{ij} = \sum_{r = 1}^p A_{ir} B_{rj}.
            \]

            Logo, a $j$-\'{e}sima coluna da matriz $C$ corresponde a
            \[
            C_{\mdot j} = \sum_{r = 1}^p B_{rj} A_{\mdot r},
            \]
            isso \'{e}, cada coluna da matriz $C$ \'{e} uma combina\c{c}\~{a}o linear das colunas de $A$.
        \end{solution}

        \part Demonstre que cada linha da matriz $C = A B$ \'{e} uma combina\c{c}\~{a}o linear das linhas de $B$. Quais s\~{a}o os coeficientes desta combina\c{c}\~{a}o linear?
        \begin{solution}
            Pela defini\c{c}\~{a}o de produto matricial temos que
            \[
            C_{ij} = \sum_{r = 1}^p A_{ir} B_{rj}.
            \]

            Logo, a $i$-\'{e}sima linhada matriz $C$ corresponde a
            \[
            C_{i \mdot} = \sum_{r = 1}^p A_{ir} B_{r \mdot}.
            \]
            isso \'{e}, cada linha da matriz $C$ \'{e} uma combina\c{c}\~{a}o linear das linhas de $B$.
        \end{solution}
    \end{parts}

    \question Demonstre que o produto entre duas matrizes triangulares superior \'{e} uma matriz triangular superior. Idem para matrizes triangular inferiores.
    \begin{solution}
        Sejam $A$, $A: n \times n$, e $B$, $B: n \times n$, duas matrizes triangulares superiores. Seja $C = A B$, ent\~{a}o
        \[
        C_{ij} = \sum_r A_{ir} B_{rj}.
        \]
        Para $1 \leq j < i \leq n$ temos que
        \[
        C_{ij} = \sum_r A_{ir} B_{rj} = 0
        \]
        pois
        \begin{itemize}
            \item se $r < j$ ent\~{a}o $A_{ir} = 0$,
            \item se $j < r < i$ ent\~{a}o $A_{ir} = B_{rj} = 0$ e
            \item se $i < r$ ent\~{a}o $B_{rj} = 0$.
        \end{itemize}
        Logo $C$ \'{e} triangular superior.

        Agora sejam $A$, $A: n \times n$, e $B$, $B: n \times n$, duas matrizes triangulares inferiores. Seja $D = A B$, ent\~{a}o
        \[
        C_{ij} = \sum_r A_{ir} B_{rj}.
        \]
        Para $1 \leq i < j \leq n$ temos que
        \[
        C_{ij} = \sum_r A_{ir} B_{rj} = 0
        \]
        pois
        \begin{itemize}
            \item se $r < i$ ent\~{a}o $B_{rj} = 0$,
            \item se $i < r < j$ ent\~{a}o $A_{ir} = B_{rj} = 0$ e
            \item se $j < r$ ent\~{a}o $A_{ir} = 0$.
        \end{itemize}
        Logo $C$ \'{e} triangular inferior.
    \end{solution}

    \question[Exerc\'{i}cio 1.1 do Trefethen\nocite{Trefethen:1997:numerical}] Seja $B : 4 \times 4$. Sobre esta matriz s\~{a}o aplicadas as seguintes opera\c{c}\~{o}es:
    \begin{enumerate}
        \item coluna 1 multiplicada por $2$;
        \item linha 3 multiplicada por $0.5$;
        \item acicionar linha $3$ \`{a} linha $1$;
        \item permutar colunas $1$ e $4$;
        \item subtrair a linha $2$ de cada uma das outras linhas;
        \item permutar linhas $4$ e $3$;
        \item deletar a coluna $1$ (tal que a dimens\~{a}o da matriz fique $4 \times 3$).
    \end{enumerate}
    \begin{parts}
        \part Escreva esta sequ\^{e}ncia de opera\c{c}\~{o}es como um produto entre $8$ matrizes. Explique cada matriz.
        \begin{solution}
            \begin{enumerate}
                \item $B M^{(1)}$, onde
                    \[
                    M^{(1)} = \begin{bmatrix}
                        2 & 0 & 0 & 0 \\
                        0 & 1 & 0 & 0 \\
                        0 & 0 & 1 & 0 \\
                        0 & 0 & 0 & 1
                    \end{bmatrix} .
                    \]
                \item $M^{(2)} B M^{(1)}$, onde
                    \[
                    M^{(2)} = \begin{bmatrix}
                        1 & 0 & 0 & 0 \\
                        0 & 1 & 0 & 0 \\
                        0 & 0 & 3 & 0 \\
                        0 & 0 & 0 & 1
                    \end{bmatrix} .
                    \]
                \item $M^{(3)} M^{(2)} B M^{(1)}$, onde
                    \[
                    M^{(3)} = \begin{bmatrix}
                        1 & 0 & 1 & 0 \\
                        0 & 1 & 0 & 0 \\
                        0 & 0 & 1 & 0 \\
                        0 & 0 & 0 & 1
                    \end{bmatrix} .
                    \]
                \item $M^{(3)} M^{(2)} B M^{(1)} M^{(4)}$, onde
                    \[
                    M^{(4)} = \begin{bmatrix}
                        0 & 0 & 0 & 1 \\
                        0 & 1 & 0 & 0 \\
                        0 & 0 & 1 & 0 \\
                        1 & 0 & 0 & 0
                    \end{bmatrix} .
                    \]
                \item $M^{(5)} M^{(3)} M^{(2)} B M^{(1)} M^{(4)}$, onde
                    \[
                    M^{(5)} = \begin{bmatrix}
                        1 & -1 & 0 & 0 \\
                        0 & 1 & 0 & 0 \\
                        0 & -1 & 1 & 0 \\
                        0 & -1 & 0 & 1
                    \end{bmatrix} .
                    \]
                \item $M^{(6)} M^{(5)} M^{(3)} M^{(2)} B M^{(1)} M^{(4)}$, onde
                    \[
                    M^{(6)} = \begin{bmatrix}
                        1 & 0 & 0 & 0 \\
                        0 & 1 & 0 & 0 \\
                        0 & 0 & 0 & 1 \\
                        0 & 0 & 1 & 0
                    \end{bmatrix} .
                    \]
                \item $M^{(6)} M^{(5)} M^{(3)} M^{(2)} B M^{(1)} M^{(4)} M^{(7)}$, onde
                    \[
                    M^{(7)} = \begin{bmatrix}
                        0 & 0 & 0 \\
                        1 & 0 & 0 \\
                        0 & 1 & 0 \\
                        0 & 0 & 1
                    \end{bmatrix} .
                    \]
            \end{enumerate}
        \end{solution}

        \part \'{E} poss\'{i}vel representar esta sequ\^{e}ncia de opera\c{c}\~{o}es como um produto entre tr\^{e}s matrizes: $A B C$ (onde $B$ \'{e} a matriz original)?
        \begin{solution}
            \'{E} poss\'{i}vel representar esta sequ\^{e}ncia de opera\c{c}\~{o}es como um produto entre tr\^{e}s matrizes: $A B C$, onde $B$ \'{e} a matriz original, $A = M^{(6)} M^{(5)} M^{(3)} M^{(2)}$ e $B = M^{(1)} M^{(4)} M^{(7)}$.
        \end{solution}
    \end{parts}

    \question Se $e_j$ denota a coluna $j$ da matriz identidade de ordem $n$ e se $A$ \'{e} uma matriz $n \times n$, qual o resultado dos produtos:
    \begin{parts}
        \part $A e_j$;
        \begin{solution}
            O resultado \'{e} a $j$-\'{e}sima coluna da matriz $A$.
        \end{solution}

        \part $e_j^t A$;
        \begin{solution}
            O resultado \'{e} a $j$-\'{e}sima linha da matriz $A$.
        \end{solution}

        \part $e_j^T A e_j$.
        \begin{solution}
            O resultado \'{e} o elemento $a_{jj}$.
        \end{solution}
    \end{parts}

    \question[Exerc\'{i}cio 3.6.4, p\'{a}gina 113, do Meyer\nocite{Meyer:2000:matrix}] Para cada matriz $A : m \times n$, demonstre que o produto $A^* A$ e $A A^*$ s\~{a}o matrizes Hermetianas. E, $A A^t$ e $A^t A$ s\~{a}o sim\'{e}tricas.
    \begin{solution}
        Pela defini\c{c}\~{a}o do produto matricial verifica-se que
        \[
        \left( A A^* \right)_{ij} = \sum_{r = 1}^n a_{ir} \cdot \overline{a_{jr}}.
        \]
        Logo,
        \[
        \left( A A^* \right)_{ji} = \sum_{r = 1}^n a_{jr} \cdot \overline{a_{ir}}
        \]
        e
        \[
        \overline{\left( A A^* \right)_{ji}} = \sum_{r = 1}^n \overline{a_{jr}} \cdot \overline{\overline{a_{ir}}} = \sum_{r = 1}^n a_{ir} \cdot \overline{a_{jr}}.
        \]
        Portanto, $A A^*$ \'{e} uma matriz Hermetiana. Para $A^* A$ o processo \'{e} an\'{a}logo e concluimos tamb\'{e}m que $A^* A$ \'{e} Hermetiana.

        Pela defini\c{c}\~{a}o do produto matricial verifica-se que
        \[
        \left(A A^t\right)_{ij} = \sum_{r = 1}^n a_{ir} a_{jr}.
        \]
        Logo,
        \[
        \left(A A^t\right)_{ji} = \sum_{r = 1}^n a_{jr} a_{ir}
        \]
        e portanto $A A^t$ \'{e} sim\'{e}trica. Para $A^t A$ o processo \'{e} an\'{a}logo e concluimos tamb\'{e}m que $A^t A$ \'{e} sim\'{e}trico.
    \end{solution}

    \question Julque verdadeiro ou falso a afirma\c{c}\~{a}o abaixo.
    \begin{quote}
        O produto entre duas matrizes sim\'{e}tricas resulta em uma matriz sim\'{e}trica.
    \end{quote}

    \begin{solution}
        Seja $A : n \times n$ e $B : n \times n$ duas matrizes sim\'{e}tricas. Ent\~{a}o $C = A B$ \'{e} tal que
        \[
        c_{ij} = \sum_{r = 1}^n a_{ir} b_{rj}
        \]
        e
        \[
        c_{ji} = \sum_{r = 1}^n a_{jr} b_{ri}.
        \]

        Para $i \neq j$, verifica-se que $c_{ij} \neq c_{ji}$ e portanto $C$ n\~{a}o \'{e} sim\'{e}trica.
    \end{solution}

    \question O tra\c{c}o de uma matriz quadrada, denotado por $\mbox{tr}(A)$, \'{e} definido como a soma das entradas da diagonal, isso \'{e}, $\mbox{tr}(A) = \sum_i a_{ii}$.
    \begin{parts}
        \part $\mbox{tr}(A) = \mbox{tr}(A^t)$?
        \begin{solution}
            Pela defini\c{c}\~{a}o de tra\c{c}o de uma matriz quadrada, temos
            \[
            \mbox{tr}(A) = \sum_i a_{ii}
            \]
            e
            \[
            \mbox{tr}(A^t) = \sum_i a_{ii}.
            \]
            Logo, $\mbox{tr}(A) = \mbox{tr}(A^t)$.
        \end{solution}

        \part $\mbox{tr}(A) = \mbox{tr}(A^*)$?
        \begin{solution}
            Pela defini\c{c}\~{a}o de tra\c{c}o de uma matriz quadrada, temos
            \[
            \mbox{tr}(A) = \sum_i a_{ii}
            \]
            e
            \[
            \mbox{tr}(A^*) = \sum_i \overline{a_{ii}}.
            \]
            Logo, $\mbox{tr}(A) \neq \mbox{tr}(A^t)$.
        \end{solution}

        \part Para matrizes $A : m \times n$ e $B : n \times m$, demonstre que $\mbox{tr}(AB) = \mbox{tr}(BA)$.
        \begin{solution}
            Pela defini\c{c}\~{a}o de produto matricial, temos
            \[
            (A B)_{ii} = \sum_{r = 1}^n a_{ir} b_{ri}
            \]
            e
            \[
            (B A)_{ii} = \sum_{r = 1}^m b_{ir} a_{ri}.
            \]

            Pela defini\c{c}\~{a}o de tra\c{c}o de uma matriz quadrada, temos
            \[
            \mbox{tr}(A B) = \sum_{s = 1}^m \sum_{r = 1}^n a_{sr} b_{rs}
            \]
            e
            \[
            \mbox{tr}(B A) = \sum_{s = 1}^n \sum_{r = 1}^m b_{sr} a_{rs} = \sum_{s = 1}^m \sum_{r = 1}^n a_{sr} b_{rs}.
            \]
            Logo, $\mbox{tr}(A B) = \mbox(B A)$.
        \end{solution}
    \end{parts}

    \question[Exerc\'{i}cio 3.6.11, p\'{a}gina 114, do Meyer\nocite{Meyer:2000:matrix}] Usando os resultados do item anterior, demonstre cada uma das afirma\c{c}\~{o}es abaixo. (As matrizes $A$, $B$ e $C$ n\~{a}o s\~{a}o necessariamente quadradas, por\'{e}m as dimens\~{o}es s\~{a}o tais que os produtos apresentados podem ser realizadas e resultam em matrizes quadradas.)

    \begin{parts}
        \part $\mbox{tr}(A B C) = \mbox{tr}(B C A) = \mbox{tr}(C A B)$.
        \begin{solution}
            Primeiro vamos mostrar que $\mbox{tr}(A B C) = \mbox{tr}(B C A)$.
            \begin{align*}
                \mbox{tr}\left( A B C \right) &= \mbox{tr}\left( A \left( B C \right) \right) \\
                &= \mbox{tr}\left( \left( B C \right) A \right) && \mbox{Exerc\'{i}cio anterior.} \\
                &= \mbox{tr}\left( B C A \right).
            \end{align*}

            Agora vamos mostrar que $\mbox{tr}\left( A B C \right) = \mbox{tr}\left( C A B \right)$.
            \begin{align*}
                \mbox{tr}\left( A B C \right) &= \mbox{tr}\left( \left( A B \right) C \right) \\
                &= \mbox{tr}\left( C \left( A B \right) \right) && \mbox{Exerc\'{i}cio anterior.} \\
                &= \mbox{tr}\left( C A B \right)
            \end{align*}

            Portanto, $\mbox{tr}(A B C) = \mbox{tr}(B C A) = \mbox{tr}(C A B)$.
        \end{solution}

        \part $\mbox{tr}(A B C)$ pode ser diferente de $\mbox{tr}(B A C)$.
        \begin{solution}
            Para mostrar que $\mbox{tr}\left( A B C \right)$ pode ser diferente de $\mbox{tr}\left( B A C \right)$ vamos apresentar um exemplo em que isso \'{e} observado.

            Sejam
            \[
            A = \begin{bmatrix}
                1 & -1 \\
                2 & 1
            \end{bmatrix}, \ B = \begin{bmatrix}
                1 & 3 \\
                -4 & 1
            \end{bmatrix}, \ C = \begin{bmatrix}
                1 & 2 \\
                -3 & 1
            \end{bmatrix}.
            \]
            Ent\~{a}o
            \begin{align*}
                A B C &= \begin{bmatrix}
                    5 & 7 \\
                    4 & -7
                \end{bmatrix}, \\
                \mbox{tr}\left( A B C \right) &= -2, \\
                B A C &= \begin{bmatrix}
                    6 & 9 \\
                    3 & -6
                \end{bmatrix}, \\
                \mbox{tr}\left( B A C \right) &= 0.
            \end{align*}
            Logo, $\mbox{tr}\left( A B C \right) \neq \mbox{tr}\left( B A C \right)$.
        \end{solution}

        \part $\mbox{tr}(A^t B) = \mbox{tr}(A B^t)$.
        \begin{solution}
            Primeiro vamos calcular $\mbox{tr}\left( A^t B \right)$:
            \begin{align*}
                \mbox{tr}\left( A^t B \right) &= \mbox{tr}(C) && C_{ij} = \sum_{r = 1}^n A_{ri} B_{rj} \\
                &= \sum_{s = 1}^n C_{ss} \\
                &= \sum_{s = 1}^n \sum_{r = 1}^n A_{rs} B_{rs}.
            \end{align*}

            Agora vamos calcular $\mbox{tr}\left( A B^t \right)$:
            \begin{align*}
                \mbox{tr}\left( A B^t \right) &= \mbox{tr}(D) && D_{ij} = \sum_{r = 1}^n A_{ir} B_{jr} \\
                &= \sum_{s = 1}^n D_{ss} \\
                &= \sum_{s = 1}^n \sum_{r = 1}^n A_{sr} B_{sr}.
            \end{align*}

            Pelas express\~{o}es acima concluimos que $\mbox{tr}\left( A^t B \right) = \mbox{tr}\left( A B^t \right)$.
        \end{solution}
    \end{parts}

    \question[Exerc\'{i}cio 3.6.12, p\'{a}gina 114, do Meyer\nocite{Meyer:2000:matrix}] Demonstre que $\mbox{tr}(A^t A) = 0$ se e somente se $A$ \'{e} matriz nula.
    \begin{solution}
        Pela defini\c{c}\~{a}o de produto matricial, temos
        \[
        (A^t A)_{ij} = \sum_{r = 1}^n a_{jr} a_{rj}.
        \]
        E pela defini\c{c}\~{a}o de tra\c{c}o de uma matriz quadrada, temos
        \[
        \mbox{tr}(A^t A) = \sum_{s = 1}^m \sum_{r = 1}^n a_{rs} a_{rs}.
        \]

        Logo, observa-se que $\mbox{tr}(A^t A) = 0$ se e somente se $A$ \'{e} matriz nula.
    \end{solution}

    \question Para cada matriz $A : n \times n$ use a fun\c{c}\~{a}o tra\c{c}o para justificar que a equa\c{c}\~{a}o matricial $A X - X A = I$ n\~{a}o admite solu\c{c}\~{a}o.
    \begin{solution}
        Dada a equa\c{c}\~{a}o matricial $A X - X A = I$, verifica-se que $\mbox{tr}(A X - X A) = \mbox{tr}(I)$.

        Pela defini\c{c}\~{a}o de tra\c{c}o de uma matriz quadrada temos que
        \[
        \mbox{tr}\left( I \right) = n
        \]
        e da defini\c{c}\~{a}o tamb\'{e}m decorre a seguinte propriedade
        \[
        \mbox{tr}(\alpha A + \beta B) = \alpha \mbox{tr}(A) + \beta \mbox{tr}(B).
        \]
        Desta propriedade temos que
        \[
        \mbox{tr}(A X - X A) = \mbox{tr}(A X) - \mbox{tr}(X A).
        \]
        Como demonstrado anteriormente, verifica-se $\mbox{tr}(A B) = \mbox{tr}(B A)$ e portanto
        \[
        \mbox{tr}(A X - X A) = \mbox{tr}\left( A X \right) - \mbox{tr}\left( X A \right) = \mbox{tr}\left( A X \right) - \mbox{tr}\left( A X \right) = 0.
        \]

        Logo, verifica-se que $\mbox{tr}(A X - X A) \neq \mbox{tr}(I)$ e conseguentemente a equa\c{c}\~{a}o matricial n\~{a}o admite solu\c{c}\~{a}o.
    \end{solution}
\end{questions}

\bibliographystyle{plain}
\bibliography{bibliography}
\end{document}
